%----------------------------------------------------------------------------------------
%	                            PORTADA
%----------------------------------------------------------------------------------------
% \makeportada{}{Daniel Carbajales Soto}{ 1° CFGS Automatización y Robótica Industrial}{}

%----------------------------------------------------------------------------------------
%                                    HEADER 
%----------------------------------------------------------------------------------------
% \header{}{Daniel Carbajales Soto}{1° CFGSARI}
%----------------------------------------------------------------------------------------
%                               TABLE OF CONTENTS
%----------------------------------------------------------------------------------------
% \maketoc
%----------------------------------------------------------------------------------------
%	                            TABLES
%----------------------------------------------------------------------------------------
% this is just a demo in case i need to make a table
%
%
%
% \begin{tabularx}{\linewidth}{|C{2cm}|Y|Y|C{2cm}|}
% \hline
% \makecell{Header 1} & \makecell{Header 2 \\ with linebreak} & \makecell{Header 3} & \makecell{Header 4} \\
% \hline
% Row 1 & This is a long cell text that automatically wraps in the Y column & Another long cell & Short \\
% \hline
% Row 2 & \multirow{2}{*}{Multirow text} & More text & Short \\
% \cline{1-1} \cline{3-4}
% Row 3 & & Extra text & Short \\
% \hline
% \end{tabularx}


%----------------------------------------------------------------------------------------
%	                            DOCUMENT
%----------------------------------------------------------------------------------------


\fchapter{2}{Ej 4 Pág. 17.}

\fsection{\boldit{\textcolor{waveBlue2}{Expresión Oral} Tarea. }}En el blog La Ecuación Digital , su autor nos habla de la convergencia IT-OT como reto clave para 
alcanzar el éxito en la digitalización industrial.
\vspace{1em}

En la noticia recoge una afirmación de Daniel Seseña , director de industria 4.0 en Minsait , que dice así : 
<< Este nuevo escenario requiere que las organizaciones dispongan de una hoja 
de ruta adaptada a su situación y que permita , de manera estructurada y gradual , adecuar la velocidad del cambio a las necesidades operativas de su producción , 
su capacidad y recursos . Además de una visión global de todos los actores , sistemas y activos de la empresa para definir una arquitectura
tecnológica que responda de manera eficiente , coherente y escalable a las necesidades de la operación >>.
\vspace{1em}

Cada estudiante o grupo de estudiantes deberá parafrasear la afirmación de Daniel Seseña , de modo que se obtenga un texto más comprensible . 
\vspace{1em}

A continuación , se pondrán en común las conclusiones alcanzadas y se explicará el significado de palabras o expresiones como << hoja de ruta >> ,
<< actores , sistemas y activos de la empresa >> y << arquitectura tecnológica escalable >> .

\fsubsection{$\boxed{\text{\boldit{Definiciones:}}}$}
\begin{itemize}[label=\textbullet { -}]
	\item \boldit{hoja de ruta:} \\
		Plan estratégico que guía la digitalización industrial de forma gradual, ajustando la frecuendia  de los cambios a las capacidades y recursos de la empresa  para una convergencia efectiva IT-OT.
	\item \boldit{actores:} \\
		Personas y equipos involucrados en procesos industriales (ej. Directores de Información, departamentos), cuya visión global fomenta la integración y la agilidad  en la cadena de valor de Industria 4.0.
	\item \boldit{sistemas:}\\
		Soluciones tecnológicas IT (ERP, CRM para datos) y OT (SCADA, PLC para control en tiempo real); su conjunto reduce costes (hasta 20\%) y mejora la productividad (3-5\%) mediante el analisis de datos 
		avanzado o ``Big Data''.
	\item \boldit{activos de la empresa:}\\
		Recursos clave como maquinaria, robots, software, propiedades y socios estratégicos.
	\item \boldit{arquitectura:}\\
		Estructura que organiza sistemas, datos y procesos, alineando ``IT'' y ``OT'' para eliminar barreras entre ellos, adoptar tecnologías nuevas y asegurar  la comunicación entre 
		los departamentos de ciberseguridad y operaciones.
	\item \boldit{tecnológica:}\\
		Herramientas digitales (IoT, IA, big data) que conectan IT y OT, transformando producción tradicional en operaciones flexibles, personalizadas y eficientes en Industria 4.0.
	\item \boldit{escalable:}\\
		Capacida de los sistemas para crecer adaptándose a demandas mayores sin perder rendimiento, permitiendo innovación continua y competitividad global.
\end{itemize}




\clearpage
%%%%%%%%%%%%%%%%%%%%%%%%%%%%%%%%%%%%%%%%%%%%%%%%%%%%%%%%%%
%%%%%%%%%%%%%%%%%%%% bibliography %%%%%%%%%%%%%%%%%%%%%%%%
% \nocite{*}                                 %%%%%%%%%%%%%
% \printbibliography                         %%%%%%%%%%%%%
%%%%%%%%%%%%%%%%%%%%%%%%%%%%%%%%%%%%%%%%%%%%%%%%%%%%%%%%%%
%%%%%% Last page in case there is no bibliography %%%%%%%%
% \label{Last Page}                           %%%%%%%%%%%%%
%%%%%%%%%%%%%%%%%%%%%%%%%%%%%%%%%%%%%%%%%%%%%%%%%%%%%%%%%%

